\documentclass[10pt,a4paper]{amsart}
\usepackage[margin=19mm]{geometry}
\usepackage{amsmath,amssymb,mathtools}
\usepackage[scr=rsfs,cal=euler]{mathalfa}
\usepackage{xfrac}
\usepackage[unicode,hidelinks,bookmarks=false]{hyperref}
\newcommand{\ie}{i.e.\ }
\newcommand{\cf}{cf.\ }
\newcommand{\resp}{respectively\ }
\newcommand{\Subsection}{\subsection}
\providecommand{\nobreakdash}{\nobreak\hbox{-}}
\setlength{\emergencystretch}{2em}
\multlinegap0pt
\usepackage{bm}
\usepackage{bbm}
\usepackage{dsfont}
\usepackage{xspace}
\usepackage{cancel}
\usepackage{mathtools}
\usepackage[shortlabels]{enumitem}
\usepackage{amsthm}
\makeatletter
\@ifundefined{equation}{\newtheorem{equation}{Equation}}{}
\@ifundefined{prop}{\newtheorem{prop}[equation]{Proposition}}{}
\@ifundefined{theorem}{\newtheorem{theorem}[equation]{Theorem}}{}
\makeatother
\DeclareMathOperator{\Hom}{Hom}
\DeclareMathOperator{\res}{res}
\newcommand{\Jc}{{\mathcal J}}
\newcommand{\mF}{{\mathbb F}}
\newcommand{\mN}{{\mathbb N}}
\newcommand{\td}[1]{\langle #1\rangle}
\newcommand{\xra}{\xrightarrow}
\renewcommand{\to}{\longrightarrow}
\title[Representation-graded Bredon homology of elementary abelian ]{Representation-graded Bredon homology of elementary abelian 2-groups}
\author{Markus Hausmann, Stefan Schwede}
\date{Source: arXiv:2403.05355v2 (CC BY 4.0)}
\begin{document}
\maketitle

\noindent\emph{Source and licence.} Representation-graded Bredon homology of elementary abelian 2-groups. Version arXiv:2403.05355v2 (CC BY 4.0), distributed under the Creative Commons Attribution 4.0 International licence, as recorded in the arXiv licence metadata for this exact version and shown on the abstract page https://arxiv.org/abs/2403.05355v2. Full rights notice: licenses/arxiv-2403.05355-CC-BY-4.0.txt.\par\smallskip

\noindent\textit{Editorial scope.} Counted unit: the Theorem whose source label is \texttt{thm:Bredon\_RO-graded} in the article (source0.tex lines 704--723, source label \texttt{thm:Bredon\_RO-graded}), delivered together with the article's own complete original proof of it (source0.tex lines 724--909, one \texttt{proof} environment of 186 lines). No mathematics is written, repaired or re-derived: statement and proof are the article's own text line for line. Required context retained uncounted: prop:x-relation (lines 423--432); prop:fundamental (lines 494--512). Every definition, display and label of those lines is retained unchanged. Omitted from this package and not redistributed: 1--422, 433--493, 513--703, 910--1279. Nothing inside a retained excerpt is omitted or altered: every retained body is delivered exactly as the article's lines are written, so no omission receipt of any kind accompanies this package. Citations: the retained excerpts carry no citation command at all, so no bibliography entry of the article is transcribed and no reference is redistributed. Reference adaptation: none was needed and none was made; every reference inside the delivered unit resolves inside it, so no unresolved reference or citation is produced. Printed slips kept literal: no printed word, symbol, hypothesis or number of the counted statement or of its proof was altered. Layout and class: the article's own class is \texttt{amsart} with packages including enumerate, amssymb, hyperref; the wrapper is the lane's pinned \texttt{amsart} editorial wrapper at 10pt on A4, which restates the theorem-like environments the retained text needs and adds the article's own macro definitions that the retained text needs (\texttt{\textbackslash Hom}, \texttt{\textbackslash Jc}, \texttt{\textbackslash mF}, \texttt{\textbackslash mN}, \texttt{\textbackslash res}, \texttt{\textbackslash td}, \texttt{\textbackslash to}, \texttt{\textbackslash xra}), with extra wrapper packages (bm, bbm, dsfont, xspace, cancel, mathtools, enumitem). The delivered numbering is the unit's own. Assets: no retained span contains a figure environment, a graphics-inclusion command, a PDF-page inclusion, a table or a TikZ picture; no image file is copied into this package, the package declares no asset dependency and no image file is redistributed with it. Licence: version arXiv:2403.05355v2 of this article is distributed under the Creative Commons Attribution 4.0 International licence, as recorded in the arXiv licence metadata for this exact version and shown on the abstract page https://arxiv.org/abs/2403.05355v2. A full rights notice is delivered at licenses/arxiv-2403.05355-CC-BY-4.0.txt. Characters: every retained excerpt is pure ASCII, so the delivered engine, XeLaTeX with its default Latin Modern text font, reports no missing character.
\begin{prop}\label{prop:x-relation}
  Let $G$ be a finite group.
  \begin{enumerate}[\em (i)]
  \item If $V$ is an orientable $G$-representation, then $\beta(t_V)=0$.
  \item Let $T$ be a set of $G$-characters whose product is 1. Then
    \[
      \sum_{\lambda\in T} a_\lambda \cdot \bigg( \prod_{\mu\in T\setminus\{\lambda\}} t_\mu\bigg)
      \ = \ 0\ .\]
\end{enumerate}
\end{prop}

\begin{prop}\label{prop:fundamental}
  Let $A$ be an elementary abelian 2-group and let $W$ be an $A$-representation
  with trivial fixed points.
  \begin{enumerate}[\em (i)]
\item For every subgroup $B$ of $A$, the restriction homomorphism
  \[ \res^A_B\ : \ H_m(A,W)\ \to \ H_{m-k}(B,W_B) \]
  is surjective, where $k=\dim(W^B)$.
\item 
  Suppose that $W=V\oplus\lambda$ for an $A$-representation $V$ and
  a nontrivial $A$-character $\lambda$ with kernel $K$.
  Then for $k=\dim(V^K)$, the following sequence is exact: 
  \[ 0 \ \to \ H_m(A,V)\ \xra{\ \cdot a_\lambda} \ H_m(A,V\oplus\lambda)\ \xra{\res^A_K} \
    H_{m-k-1}(K,V_K)\ \to\ 0  \]

  \item The  $\mF_2$-vector space $H_m(A,W)$ is spanned by the classes $a_U\cdot t_V$
  for all $A$-representations $U$ and $V$ such that $U\oplus V= W$
  and $m=\dim(V)$.
\end{enumerate}
\end{prop}

\begin{theorem}\label{thm:Bredon_RO-graded}
  Let $A$ be an elementary abelian 2-group.
  \begin{enumerate}[\em (i)]
  \item
    The representation-graded Bredon homology ring $H(A,\star)$
    is generated as an $\mF_2$-algebra
    by the classes $a_\lambda$ and $t_\lambda$ for all nontrivial $A$-characters $\lambda$.
\item 
  The kernel of the surjective homomorphism of $\mF_2$-algebras
  \[ \epsilon_A \ : \  \mF_2[a_\mu,t_\mu\colon \mu\in A^\circ] \ \to \ H(A,\star) \]
  is the ideal generated by the polynomials
  \[  r(T)\ = \
    \sum_{\lambda\in T}  a_\lambda \cdot\bigg( \prod_{\mu\in T\setminus\{\lambda\}}\, t_\mu\bigg) \]
  for all minimally dependent subsets $T$ of $A^\circ$.
\item
  Every set of homogeneous elements that generates the kernel
  of $\epsilon_A$ contains the polynomials $r(T)$
 for all minimally dependent subsets $T$ of $A^\circ$.
\end{enumerate}
\end{theorem}

\begin{proof}
  Part (i) was shown in Proposition \ref{prop:fundamental} (iii),
  and is repeated here for easier reference.

  (ii)
  For the course of the proof we write $I(A)$ for the ideal of the polynomial ring 
  $\mF[a_\mu,t_\mu\colon \mu\in A^\circ]$ generated by the polynomials $r(T)$
  for all minimally dependent subsets $T$ of $A^\circ$. 
  Since minimally dependent sets of characters multiply to 1,
  Proposition \ref{prop:x-relation} (ii) shows that $I(A)\subseteq \ker(\epsilon_A)$;
  so it remains to show the reverse inclusion.

  We argue by induction on the rank of $A$.
  The induction starts with the trivial group, where there is nothing to show.
  Now we let $A$ be a nontrivial elementary abelian 2-group, and we assume part (ii)
  for all proper subgroups of $A$.
  In the inductive step, we shall make use of the polynomials
  \[  r(T)\ = \
    \sum_{\lambda\in T}  a_\lambda \cdot\bigg( \prod_{\mu\in T\setminus\{\lambda\}}\, t_\mu\bigg) \]
  for arbitrary subsets $T$ of $A^\circ$, not necessarily minimally dependent;
  we alert the reader that if the elements of $T$ do not multiply to the trivial character,
  then the polynomial $r(T)$ will {\em not} map to 0 under $\epsilon_A$.
  
  We let $\lambda$ be a nontrivial $A$-character, with kernel $K$.
  We let $\Jc$ denote the homogeneous ideal in the ring
  $\mF_2[a_\mu,t_\mu\colon  \mu\in A^\circ\setminus\{\lambda\}]$ consisting of those
  elements $y$ such that $\res^A_K(\epsilon_A(y))=0$.
  We emphasize that elements of $\Jc$ are polynomials that do not involve the variables
  $a_\lambda$ nor $t_\lambda$.
  We shall prove two properties of this ideal:\smallskip

  (a) The ideal $\Jc$ is generated by the polynomials $r(T)$ for all
  subsets $T$ of $A^\circ\setminus\{\lambda\}$
  such that either $T$ is minimally dependent, or $T\cup\{\lambda\}$
  is minimally dependent.\smallskip

  (b) In the graded ring   $\mF_2[a_\mu,t_\mu\colon  \mu\in A^\circ]$,
  the relation $t_\lambda\cdot \Jc\subset \td{I(A),a_\lambda}$ holds.\smallskip

  Proof of (a). We write $J_{A\backslash\lambda}$ for the submonoid of $J_A$ consisting of
  the classes of $A$-representations $V$ that do not involve the character $\lambda$.
  Equivalently, $J_{A\backslash\lambda}$ contains the $A$-representations
  with trivial $A$-fixed points such that $V^K=0$.
  The restriction homomorphism $\res^A_K:\mN\times J_{A\backslash\lambda}\to\mN\times J_K$
  lets us inflate $(\mN\times J_K)$-graded rings $R$
  to $(\mN\times J_{A\backslash\lambda})$-graded rings $(\res^A_K)^*(R)$ by setting
  \[ (\res^A_K)^*(R)(k,V)\ = \ R(k,V|_K)\ . \]
  We alert the reader that a grading-inflated polynomial algebra
  is no longer a polynomial algebra.
  With this grading convention, a morphism
  of $(\mN\times J_{A\backslash\lambda})$-graded $\mF_2$-algebras
  \[  \rho^A_K\ : \   \mF_2[a_\mu,t_\mu\colon  \mu\in A^\circ\setminus\{\lambda\}] \ \to \
  (\res^A_K)^*\left(  \mF_2[a_\nu,t_\nu\colon  \nu\in K^\circ]  \right) \]
  is given by sending $a_\mu$ to $a_{\mu|K}$ and $t_\mu$ to $t_{\mu|K}$. 
  Moreover, the ideal $\Jc$ is precisely the kernel of the composite homomorphism
  \begin{align*}
    \mF_2[a_\mu,t_\mu\colon  \mu\in A^\circ\setminus\{\lambda\}] \
    &\xra{\ \rho^A_K\ } \
      (\res^A_K)^*\left(  \mF_2[a_\nu,t_\nu\colon  \nu\in K^\circ]  \right)\\
    &\xra{(\res^A_K)^*(\epsilon_K)} (\res^A_K)^* (  H(K,\star) ) \ .
  \end{align*}
  The kernel of $\rho^A_K$ is the ideal generated by the homogeneous polynomials
  \[
    a_\mu t_{\mu\lambda} + t_\mu a_{\mu\lambda}\ = \ r(\{\mu,\mu\lambda\})
  \]
  for all $\mu\in A^\circ\setminus\{\lambda\}$.
  The set $\{\mu,\mu\lambda\}$ is independent
  and $\{\mu,\mu\lambda\}\cup\{\lambda\}$ is minimally dependent.
  So the polynomials $r(\{\mu,\mu\lambda\})$ are among those of the second kind specified in (a).

  By the inductive hypothesis for the subgroup $K$,
  the kernel of $\epsilon_K: \mF_2[a_\nu,t_\nu\colon  \nu\in K^\circ] \to H(K,\star)$
  is generated by the polynomials $r(S)$ for all minimally dependent subsets $S$
  of $K^\circ$. So the kernel of the homomorphism
  \[  (\res^A_K)^*(\epsilon_K)\ : \ 
    (\res^A_K)^*\left(  \mF_2[a_\nu,t_\nu\colon  \nu\in K^\circ]  \right)
    \ \to \  (\res^A_K)^* (  H(K,\star) ) \]
  is the ideal generated by the same polynomials $r(S)$, but each occurring
  multiple times in different degrees, namely for all subsets $T$ of
  $A^\circ\setminus\{\lambda\}$ such that 
  \[ {\bigoplus}_{\mu\in T}\ \mu|_K \ = \ {\bigoplus}_{\nu\in S}\ \nu\ . \]
  This condition means that for each character $\nu\in S$, the set $T$
  contains exactly one of the two extensions of $\nu$ to an $A$-character.
  Because $S$ is a minimally dependent subset of $K^\circ$, the $A$-character
  \[ {\prod}_{\mu\in T}\ \mu \quad \in \ \Hom(A,C)\]
  then restricts to the trivial character on the subgroup $K$.
  Hence this product is either 1 or $\lambda$.
  In the first case, $T$ is a minimally dependent subset of $A^\circ\setminus\{\lambda\}$.
  In the second case, $T$ is an independent subset of $A^\circ\setminus\{\lambda\}$
  such that $T\cup\{\lambda\}$ is minimally dependent.

  We have now shown that some of the polynomials $r(T)$ with $T$ as specified in (a)
  generate the kernel of $\rho^A_K$, and the images of the others under $\rho^A_K$
  generate the kernel of $(\res^A_K)^*(\epsilon_K)$. So all those polynomials
  together generate the kernel of the composite, and hence the ideal $\Jc$.
  This completes the proof of (a).\medskip

 
    Proof of (b). It suffices to show that the two types of generating polynomial for
  $\Jc$ specified in (a) have the desired property.
  If $T$ is a minimally dependent subset of $A^\circ\setminus\{\lambda\}$,
  then the polynomial $r(T)$ belongs to the ideal $I(A)$.
  Hence also  $t_\lambda\cdot r(T)\in I(A)$.
  If $T$ is a subset of $A^\circ\setminus\{\lambda\}$ such that
  $T\cup\{\lambda\}$ is minimally dependent, then the relation
  \[  t_\lambda \cdot r(T)\
    = \   r(T\cup\{\lambda\}) \ + \ a_\lambda \cdot {\prod}_{\mu\in T}\ t_\mu  \]
  holds in the ring $\mF_2[a_\mu,t_\mu\colon \mu\in A^\circ]$,
  and witnesses that the left hand side lies in the ideal $\td{I(A),a_\lambda}$.
  This completes the proof of (b).\medskip
  
  Now we complete the inductive step, showing that $\ker(\epsilon_A)\subseteq I(A)$.
  We prove this for all homogeneous pieces $H_m(A,W)$,
  where $W$ is an $A$-representation with $W^A=0$,
  by induction on the dimension of $W$.
  For $W=0$ we are considering integer degrees,
  where source and target of $\epsilon_A$ both
  consist only of a copy of $\mF_2$ in degree 0.
  
  Now we suppose that $W\ne 0$. 
  We let $f$ be a homogeneous polynomial in $\mF_2[a_\mu,t_\mu\colon  \mu\in A^\circ]$
  of degree $(k,W)$ with $\epsilon_A(f)=0$.
  We choose an $A$-character $\lambda$ and
  an $A$-representation $V$ with $W=V\oplus\lambda$.
  We let $K$ denote the kernel of $\lambda$.
  We write $f=t_\lambda\cdot y+a_\lambda\cdot z$ for some homogeneous polynomials $y,z$
  of degrees $(k-1,V)$ and $(k,V)$, respectively. Then
  \[
    \res^A_K(\epsilon_A(y))\ = \
    \res^A_K( t_\lambda\epsilon_A(y) + a_\lambda \epsilon_A(z))
    \ = \     \res^A_K( \epsilon_A(f))    \ = \ 0 \ . \]
  
  Case 1: The $A$-character $\lambda$ occurs in $W$ with multiplicity at least 2.
  Then there is an $A$-representation $U$ such that $V=U\oplus \lambda$.
  Because $\res^A_K(\epsilon_A(y))=0$, there is a class $u\in H_{k-1}(A,U)$
  such that $\epsilon_A(y)=a_\lambda\cdot u$.
  Then 
  \[ a_\lambda\cdot (t_\lambda\cdot u+\epsilon_A(z))\ =\
    t_\lambda\cdot\epsilon_A(y)+a_\lambda\cdot\epsilon_A(z)\ = \ \epsilon_A(f)\ =\ 0\ . \]
  Because multiplication by $a_\lambda$ is injective,
  we deduce that $t_\lambda\cdot u=\epsilon_A(z)$.
  We choose a homogeneous polynomial $g$ in $\mF_2[a_\mu,t_\mu\colon\mu\in A^\circ]$
  of degree $(k-1,U)$ such that $\epsilon_A(g)=u$.
  Then $y+a_\lambda g$ and $z+t_\lambda g$ lie in the kernel of $\epsilon_A$.
  Since $U\oplus\lambda=V$ has smaller dimension than $W$,
  the classes $y+a_\lambda g$ and $z+t_\lambda g$
  are in the ideal $I(A)$, by induction. So
  \[  f\ =\ t_\lambda\cdot y+a_\lambda\cdot z\ = \
    t_\lambda\cdot(y+a_\lambda g) +a_\lambda\cdot(z+t_\lambda g)  \]
  also lies in the ideal $I(A)$.\smallskip

  Case 2: The $A$-character $\lambda$ occurs in $W$ with multiplicity 1.
  Then $\lambda$ does not occur in $V$, and so the polynomial $y$
  does not involve the variable $a_\lambda$ nor $t_\lambda$.
  Because $\res^A_K(\epsilon_A(y))=0$, the polynomial $y$ belongs to the ideal $\Jc$.
  By property (b) of the ideal $\Jc$, the class $t_\lambda \cdot y$ then belongs
  to the ideal generated by $I(A)$ and $a_\lambda$.
  Hence there is a homogeneous polynomial $h$ of degree $(k,V)$
  such that $t_\lambda\cdot y$ is congruent to $a_\lambda\cdot h$ modulo the ideal $I(A)$.
  So the polynomial $f$ in the kernel of $\epsilon_A$ satisfies
  \[ f\ = \ t_\lambda\cdot y + a_\lambda\cdot z\ \equiv \ a_\lambda\cdot (h+z)  \text{\qquad modulo $I(A)$.}\]
  Because $I(A)$ is contained in the kernel of $\epsilon_A$, we deduce the relation
  \[ a_\lambda\cdot\epsilon_A(h+z) \ = \  \epsilon_A(a_\lambda\cdot(h+z)) \ = \ \epsilon_A(f) \ = \  0 \ .\]
  Because multiplication by $a_\lambda$ is injective, we conclude that $\epsilon_A(h+z)=0$.
  Since $V$ has smaller dimension than $W=V\oplus \lambda$, the class $h+z$
  lies in the ideal $I(A)$ by induction. So also $a_\lambda(h+z)$, and hence the class $f$,
  lies in the ideal $I(A)$. This finishes the inductive step, and hence the proof of part (ii).

  (iii)
  Source and target of the homomorphism $\epsilon_A$
  are graded by the abelian monoid of isomorphism classes of $A$-representations.
  This abelian monoid is free on the classes of the $A$-characters
  and thus admits a compatible partial order by declaring $V\leq W$
  if $V$ is isomorphic to a direct summand of $W$.
  Hence all products of a homogeneous polynomial $f$ with other homogeneous elements of
  $\mF_2[a_\mu,t_\mu\colon \mu\in A^\circ]$ have degrees greater or equal than that of $f$.
  By (ii), the kernel of $\epsilon_A$ is generated in gradings 
  that contain the sum of a linearly dependent set of characters.
  So the kernel of $\epsilon_A$ is nontrivial only in degrees 
  that contain the sum of a linearly dependent set of characters.
  The generating relations $r(T)$ specified in (ii) lie in degrees
  that are minimal with this property, and they are the unique nontrivial
  elements in the kernel of $\epsilon_A$ in their degrees.
  So each of the relations $r(T)$ specified in (ii) is necessary to generate
  the kernel of $\epsilon_A$.
\end{proof}

\end{document}
